% PDF settings for the single-file Bhagavad Gita rendering.
% This PDF intentionally omits the Sanskrit text. The glossary is placed at the end.

\usepackage{microtype}
\usepackage{enumitem}
\usepackage{setspace}
\usepackage{graphicx}
\usepackage{xcolor}
\usepackage{eso-pic}

% Full-page cover used before the table of contents.  The starred shipout
% hooks apply to this page only, so later pages retain their normal white
% background.  The image is fitted without cropping and centred on A5 paper.
\newcommand{\gitacover}[1]{%
  \thispagestyle{empty}%
  \AddToShipoutPictureBG*{%
    \AtPageLowerLeft{\color[HTML]{FDF7EB}\rule{\paperwidth}{\paperheight}}%
  }%
  \AddToShipoutPictureFG*{%
    \AtPageCenter{%
      \makebox(0,0){%
        \includegraphics[width=\paperwidth,height=\paperheight,keepaspectratio]{#1}%
      }%
    }%
  }%
  \null\clearpage%
}

\setstretch{1.04}
\setlength{\parindent}{0pt}
\setlength{\parskip}{0.22em}

% Keep the front-matter table of contents on page 2 only.  The slightly
% tighter leading compensates for the more generous page margins, while the
% final clear page reserves page 3 exclusively for the introductory note.
\let\gitastandardtableofcontents\tableofcontents
\renewcommand{\tableofcontents}{%
  \begingroup
  \setstretch{0.96}%
  \gitastandardtableofcontents
  \endgroup%
  \clearpage%
}

\setlist[description]{
  style=nextline,
  leftmargin=0pt,
  labelsep=0pt,
  itemsep=0.35em,
  parsep=0pt,
  topsep=0.2em
}
